\section{Introduction}

Learning-based reconstruction has become the default computational layer in imaging through scattering media. Networks trained on paired speckle--object data now recover object structure across diffuser changes \cite{li2018deep}, under dynamic scattering \cite{liu2024learning}, and in real time with modest compute \cite{zhang2024memoryless}. As reported fidelity improves, however, a separate line of work has established that these reconstructions can \emph{fabricate}: they invent spectral or semantic content that the measurement does not physically constrain \cite{zhang2026physical,tivnan2024hallucination}. Fabrication is not a rare failure; it is the expected behavior of an underdetermined inverse problem solved by a prior-rich model, and it has been analyzed mechanistically for scattering imaging \cite{zhang2026physical} and quantified by dedicated assessment frameworks \cite{long2026universal}.

What remains missing is a \emph{causal control knob}. Existing analyses explain why fabrication happens, and assessment frameworks measure how much of it a given model exhibits, but neither establishes which training-time choice---a choice available for free---causally governs it. This distinction matters because a correlational audit cannot tell a practitioner what to change: fidelity metrics, reconstruction losses, and reported ``generalization to unseen diffusers'' all co-vary with data, architecture, and optimization budget at once. One preprocessing decision sits directly on the boundary between what the measurement contains and what the network must invent: the spatial bandwidth of the \emph{training target}. Band-unlimited ground-truth targets (e.g., clean digital renderings or sharpened images) instruct the network to reproduce spatial frequencies that the measured speckle transfer function cannot carry; a bandwidth-matched target instructs it only to reproduce what the measurement can certify. The natural hypothesis---that band-unlimited targets are a major cause of measured fabrication---has, to our knowledge, never been tested causally, against a measured information support, or against a compute-matched control.

This paper tests that hypothesis in a fully specified speckle-imaging protocol and quantifies the result against two reference points: the measured MTF support of the forward operator, and a Wiener linear baseline. Our contributions are:

\begin{itemize}
\item \textbf{A causal, compute-matched intervention.} From an identical initialization and with identical additional optimization, fine-tuning against bandwidth-matched targets reduces the reconstruction's out-of-band spectral energy to $0.34\pm0.11$ (matched) and $0.20\pm0.08$ (half) of the compute-matched raw-target control across two architectures and three seeds, with 6/6 units reducing (one-sided sign test $p\approx0.016$) while the no-reference forward-consistency endpoint changes by only $+0.026$ on a $0.36$--$0.40$ baseline (Sec.~\ref{sec:main-control}).
\item \textbf{A capacity effect on fabrication.} In the two tested capacities, raw fabrication is higher for the larger network, and the suppression under bandwidth-matched fine-tuning is stronger there too (Sec.~\ref{sec:capacity}).
\item \textbf{A Wiener-anchored trade-off.} A Wiener linear baseline bounds the fabrication budget ($2$--$6\times$ below the raw-target and matched arms across a regularization/support grid, with the half-bandwidth arm approaching its range) but remains under-fitted: even with optimal scalar rescaling its relative forward error stays about $40\%$ above the learned arms, which therefore occupy a distinct operating point rather than dominating (Sec.~\ref{sec:wiener}).
\item \textbf{A negative control on the class-label explanation.} Cross-diffuser transfer loss varies smoothly with statistical distance and shows no changepoint under three distance metrics, supporting protocol-level reporting (measured bandwidth, measured distance) over class labels such as ``same diffuser class'' (Sec.~\ref{sec:negative}).
\item \textbf{Robustness and a real-data boundary condition.} The suppression effect survives a position-dependent forward operator (Sec.~\ref{sec:svpsf}); an assessment on a public experimental dataset quantifies system parameters and shows that quantitative real-data evaluation requires a dedicated forward-calibration stage (Sec.~\ref{sec:realdata}).
\end{itemize}

Our strongest single result is the paired fine-tune-to-control reduction in Fig.~\ref{fig:paired} and Fig.~\ref{fig:hero}: from the same starting weights, changing only the bandwidth of the fine-tuning target---a decision that costs nothing and changes no architecture---cuts a model's out-of-band spectral fabrication to roughly one third of the identically optimized control.

The remainder of this paper is organized as follows. Section~\ref{sec:related} positions the work. Section~\ref{sec:setup} defines the protocol and endpoints. Section~\ref{sec:experiments} presents the causal intervention, capacity trend, Wiener anchor, robustness checks, and the distance-law negative control. Section~\ref{sec:realdata} reports the real-data assessment. Section~\ref{sec:discussion} discusses limitations and concludes.
